\chapter{GRH怎样进入一个平均问题}\label{ch:characters}
\section{用特征检测一个子群是否太小}
我们要证明 $H_t$ 是模 $n=t^2+t+1$ 的全部单位群。有限交换群有一个方便的检测器：如果一个子群 $H$ 真小于群 $\mathcal U$，商群 $\mathcal U/H$ 非平凡，因而存在一个非平凡乘法同态
\[\chi:\mathcal U\longrightarrow\{z\in\CC:|z|=1\}
\]
在 $H$ 上恒等于一。将它在非单位处定义为零，就得到一个Dirichlet特征。

因此，$H_t$ 不满等价于存在非主特征满足
\[\chi(-1)=1,\qquad\chi(d)=1\quad(1\le d\le t,\ (d,n)=1).
\]
“非主”表示它不在所有单位上恒等于一；$\chi(-1)=1$ 通常称为偶特征。要排除它，只需在 $t$ 以内找到一个与 $n$ 互素的整数 $a$，使 $\chi(a)\ne1$。

这不只是二次剩余问题。二次特征仅取 $\pm1$，但一般商群可能有三阶、五阶等特征。只检查平方剩余会漏掉别的乘法障碍。非单位处的零值也不算找到非剩余；准确要求是 $\chi(a)\notin\{0,1\}$。\src{S16}\src{S30}

\section{当总位置数是素数，有一个鸽巢证明}
\begin{proposition}
若 $n=t^2+t+1$ 是素数，则 $H_t=(\Z/n\Z)^\times$。
\end{proposition}
\begin{proof}
给任意非零余数 $a$，考虑 $(t+1)^2$ 个数
\[ai+j\pmod n,\qquad0\le i,j\le t.
\]
它们多于 $n$，所以两项相同。相减得 $a(i-i')=j'-j$。两个差都不为零：若一个为零，由 $n$ 为素数及差的绝对值小于 $n$，另一个也为零，与两项不同矛盾。于是
\[a\equiv\frac{j'-j}{i-i'}\pmod n,
\qquad0<|i-i'|,|j'-j|\le t.
\]
分子分母及符号都在 $H_t$ 的生成元中。
\end{proof}

合数模数时，这个证明的分母可能不是单位。一般“平方根以内的小单位总能生成全部单位群”确实是假的。例如模105时，不超过 $\sqrt{105}$ 的单位只有 $1,2,4,8$；加上负一生成的群有24个元素，而全部单位有48个。105并不是这里的 $t^2+t+1$，所以这反驳的是泛化的证明策略，并不反驳平方族猜想。

\section{无条件文献已经说明大参数一定成功}
Pollack的2017年定理给出：任取 $\varepsilon>0$，存在统一的 $m_0(\varepsilon)$ 和 $\kappa>0$，对每个 $m>m_0$ 及每个非主特征模 $m$，在
\[m^{1/(4\sqrt e)+\varepsilon}
\]
以内有多于 $m^\kappa$ 个素数 $\ell$ 满足 $\chi(\ell)\notin\{0,1\}$。这里只需要其中一个。该定理允许合数模数及非本原特征，符合我们的输入。\cite{Pollack2017}

取 $\varepsilon=1/4-1/(4\sqrt e)>0$，界变成 $n^{1/4}$。对 $t\ge2$，有 $n^{1/4}<t$，所以每个充分大的 $t$ 都满足 $H_t$ 满群。

这已经是无条件结果。它留下的困难很具体：现有引用没有提供足够可用的明确起点 $T_0$。因此，验证 $t<1000$，再说“充分大时成立”，中间仍可能有一段没有覆盖；两句话不能直接拼成全参数定理。

\section{GRH提供一个可计算的交界点}
在Dirichlet $L$ 函数的GRH下，Bach的小非剩余界可写为：每个非主特征模 $n$，都有
\[a\le3(\log n)^2,\qquad\chi(a)\notin\{0,1\}.
\]
项目采用的是Martin--Pollack论文首页对这一界的明确陈述；Bach原论文列在参考文献中，本轮未取得并重读其全文。\cite{MartinPollack2013,Bach1990}

当 $t\ge1000$，有
\[3\bigl(\log(t^2+t+1)\bigr)^2<t.
\]
在 $t=1000$，左边小于573。若把左边记为 $f(t)$，则
\[f'(t)=6\log(t^2+t+1)\frac{2t+1}{t^2+t+1}
<\frac{24\log(t+1)}t<1\qquad(t\ge1000).
\]
所以右边 $t$ 此后增长更快，不等式持续成立。于是GRH给出全部 $t\ge1000$ 的统一覆盖。

剩余 $2\le t<1000$ 共998个参数已经有精确有限证书。其中188个使 $n$ 为素数，可以用上一节证明；其余810个模数使用完整单位群生成核验。二者与GRH范围合起来，得到
\begin{theorem}[平方族的条件性完成]
假设Dirichlet $L$ 函数的GRH。对每个 $t\ge2$，$k=t^2$、$n=t^2+t+1$ 的有限平均条件恰为 $G=1$。
\end{theorem}
这些998个小参数的结论本身是无条件的。对任何使 $t^2+t+1$ 为素数的参数，也有无条件结论，无需限制 $t$ 的大小。

\section{有限证书究竟验证什么}
模 $n$ 的单位群 $\mathcal U$ 先通过CRT分解成素数幂模数的单位群，再写成循环群的直积。每个小整数由一个指数向量表示。要判断这些向量是否生成整个群，可以对每个整除群阶的素数 $\ell$，检查它们在
\[\mathcal U/\mathcal U^\ell
\]
这个有限维 $\F_\ell$ 向量空间中是否张成。

理由很短：若生成子群真小，非平凡有限商群必有一个素数阶商；在对应的 $\mathcal U/\mathcal U^\ell$ 中就不可能满秩。反之，只要每个素数层都满秩，商群就只能平凡。

因此证书核验的是每个指定有限群的完整生成性。它不靠抽样单位，也不把近千个成功参数外推到所有 $t$。源包保存原证书、核验程序及文献访问记录。

\section{现有显式Burgess结果为何还没有直接去掉GRH}
一般短特征和估计当然与此有关，但用于全参数证明时必须同时核对模数条件、特征条件、常数和阈值。项目调查的一份2025年显式结果具有很大的适用起点，并有本原特征或无立方因子等限制。这里 $t^2+t+1$ 不总无立方因子：$t=18$ 时它等于 $343=7^3$。

该来源及其适用版本见\cite{ExplicitBurgess2025}；详细访问等级保存在来源S30中。

所以当前“去掉GRH”的准确任务，是给 $H_t$ 一个覆盖所有参数的无条件证明，或者给一个明确且可衔接的无条件阈值。平方族的物理返回不需要重做；另一方面，现有文献的渐近界也尚未完成这最后的全参数衔接。
